% Chapter 2 content merged into Chapter 1 in v170 (M1 macro cut).
% This file is now \input{} from inside ch1_introduction.tex as a
% \section, with internal \section levels demoted to \subsection
% and \subsection levels demoted to \subsubsection.

\section{Related Research}
\label{ch:background}\phantomsection\label{sec:bg_hotspot}

This section reviews prior work in viral mutation hotspot detection and identifies the benchmarking gap that motivated MutBench.

\subsection{Mutation Hotspot Detection Methods}
\label{subsec:bg_hotspot_methods}

Mutation hotspots are regions within a genome where mutations are concentrated in a non-uniform manner, likely reflecting functionally important sites under structural or immunological selective pressure.
The concept originates from the observation that mutations are not distributed uniformly across viral genomes: certain positions or regions accumulate mutations at rates far exceeding the background, driven by positive selection for immune evasion, receptor binding optimization, or adaptation to new host environments.
For RNA viruses, which exhibit mutation rates of $10^{-3}$ to $10^{-5}$ substitutions per nucleotide per replication cycle~\cite{holmes2009evolution,sanjuan2010viral}, the interplay between high mutation supply and strong selective pressure creates pronounced mutational clustering in functionally critical regions.
Identifying hotspots in viral genomes is critical for three practical applications: (1)~predicting variant emergence by identifying positions predisposed to future mutations, (2)~selecting vaccine targets by avoiding highly mutable epitopes or targeting conserved adjacent regions, and (3)~developing broadly effective therapeutics that target structurally constrained positions unlikely to mutate without fitness cost.
Existing hotspot detection approaches can be categorized into frequency-based, entropy-based, density-based, and selection-based methods.

\textbf{Frequency-based approaches} calculate mutation frequency at the amino acid or nucleotide level and designate positions with high frequency as hotspots.
While intuitive and easy to implement, they have two fundamental limitations.
First, \textbf{lineage bias}: when a specific lineage is overrepresented in the database, its defining mutations receive disproportionately high frequency scores.
Mutations such as D614G, observed in nearly all modern SARS-CoV-2 sequences, reflect genuine fitness advantages, but the rate of near-fixation is amplified by founder effects during the early pandemic, making it difficult to distinguish the relative contributions of selection and genetic drift from frequency data alone.
Second, \textbf{neglect of rare variants}: mutations that occur at low frequency but exhibit diverse forms---potentially indicating immune evasion or viral adaptation---are systematically overlooked.

\textbf{Entropy-based approaches} use Shannon entropy to measure nucleotide diversity at a given position.
Mullick et al.~\cite{mullick2021entropy} combined Shannon entropy with K-means clustering to detect hotspots in the SARS-CoV-2 Spike protein.
However, conventional nucleotide entropy reflects the distribution of all nucleotides regardless of whether mutation has occurred, failing to capture effective diversity conditioned on the occurrence of mutation.
Furthermore, K-means clustering requires the number of clusters to be specified in advance and does not account for density structure.
Large-scale frequency-based analyses have been applied to SARS-CoV-2~\cite{koyama2020variant,obermeyer2022analysis}, but these studies used single information types (mutation frequency or fitness proxies) without systematic comparison against alternative scoring approaches.

\textbf{Cancer genomics methods.}
In cancer genomics, a rich ecosystem of methods has been developed to detect somatic mutation hotspots.
MutSigCV~\cite{lawrence2013mutsigcv} identifies significantly mutated genes by modeling background mutation rates that account for gene expression, replication timing, and chromatin state.
OncodriveCLUST~\cite{tamborero2013oncodriveclust} detects genes with significant spatial clustering of mutations in protein sequence space, under the assumption that driver mutations cluster in functional domains.
HotMAPS~\cite{tokheim2016hotmaps} maps mutations onto protein three-dimensional structures to identify spatial hotspots in 3D, and MutSpot~\cite{munoz2020mutspot} models position-specific background mutation rates using epigenomic features.
While these tools represent sophisticated approaches to hotspot detection, they target somatic mutations in tumor genomes, which are more than 1,000 times larger than viral genomes and involve fundamentally different mutation dynamics and selective pressures (neutral somatic mutation followed by clonal selection vs.\ population-level viral evolution), making direct transfer infeasible.

An important limitation of frequency-based mutation analysis is \textbf{phylogenetic non-in\-de\-pen\-dence}: mutations observed in closely related sequences are not independent events~\cite{felsenstein1985phylogenies}.
Phylogeny-aware methods such as TreeTime~\cite{sagulenko2018treetime}, BEAST~\cite{suchard2018beast}, and HyPhy~\cite{kosakovsky2020hyphy} can address this but require high-quality phylogenetic trees and face scaling and placement challenges in densely sampled datasets, motivating ultra-fast tree-placement frameworks such as UShER~\cite{turakhia2020ultrafast}.
Several tools detect homoplasic (independently recurrent) mutations on phylogenies.
HomoplasyFinder~\cite{crispell2019homoplasyfinder} identifies positions where the same mutation arises independently on multiple branches of a phylogenetic tree, which provides evidence for convergent evolution driven by positive selection rather than shared ancestry.
ConDor~\cite{morel2024condor} extends this approach to detect convergent amino acid substitutions across large protein alignments, using statistical tests to distinguish genuine convergence from chance recurrence.
Homoplasy-based approaches are valuable because they explicitly account for phylogenetic non-independence, but they identify individual recurrent sites rather than performing the spatial clustering and multi-information-type comparison that MutBench provides.
In MutBench, homoplasy is incorporated as one of the 10 underlying biological information families, which are expanded into 20 scoring channels (with normalization or compositing variants) and evaluated across pathogens (Section~\ref{sec:information_analysis}).

Simpler approaches to hotspot identification include sliding window methods, which scan fixed-length windows across the genome and flag those exceeding a mutation count threshold.
While computationally straightforward, sliding window methods suffer from two inherent limitations: the window size must be specified in advance (too small misses extended hotspots; too large merges distinct hotspots), and they impose a fixed spatial resolution that cannot adapt to the variable density of mutations across different genomic regions.
Signal processing techniques such as wavelet transforms offer a multi-scale alternative, decomposing mutation frequency profiles to identify localized peaks at multiple spatial scales simultaneously.
However, these methods treat the genome as a one-dimensional signal and do not incorporate biological information about mutation importance or selective pressure.

\textbf{Selection-based methods} represent a fundamentally different approach.
Rather than detecting spatial clusters of mutations, codon-based methods estimate the ratio of non-synonymous to synonymous substitution rates ($\omega = dN/dS$) at each codon position.
The HyPhy phylogenetic platform~\cite{kosakovsky2020hyphy} packages a family of such tests: MEME (Mixed Effects Model of Evolution)~\cite{murrell2012meme} detects episodic positive selection at sites where $\omega > 1$ on a subset of phylogenetic branches, while FUBAR (Fast Unconstrained Bayesian AppRoximation)~\cite{murrell2013fubar} provides a Bayesian estimate of per-site selection pressure pooled across the whole tree.
Closely related tests in the same family include FEL (Fixed Effects Likelihood) and SLAC (Single Likelihood Ancestor Counting)~\cite{kosakovskypond2005fel}, which provide site-level $\omega$ estimates under a maximum-likelihood and a counting framework respectively, and BUSTED (Branch-Site Unrestricted Statistical Test for Episodic Diversification)~\cite{murrell2015busted}, which tests for episodic diversifying selection along a designated set of branches.
These methods provide a principled evolutionary framework for identifying positions under positive selection, but they detect individual selected sites rather than spatial clusters of hotspot regions, and they require reliable phylogenetic trees and codon-aligned sequences---prerequisites that are not always available for all pathogens.
MutBench incorporates four phylogenetic selection-pressure scoring channels---three FUBAR-derived posterior summaries (FUBAR posterior probability, FUBAR positive-selection indicator, FUBAR Bayes factor) plus a Nei--Gojobori dN/dS proxy (Chapter~\ref{ch:mutbench}, Section~\ref{sec:scoring_detection})---enabling direct comparison of selection-based features against frequency-based and PLM-based alternatives within a unified framework. FEL, SLAC, and BUSTED are reviewed here for completeness of the HyPhy family but are not benchmarked as scoring channels in the present panel; coverage of the broader selection-test family is left to future work (Chapter~\ref{ch:discussion}, Section~\ref{sec:future_work}).

\subsection{Density-Based Clustering}
\label{subsec:bg_density_clustering}

DBSCAN~\cite{ester1996dbscan} identifies density-connected clusters and noise from $\varepsilon$-neighborhoods and core points, allowing arbitrary cluster shapes without a pre-set cluster count.
Its fixed global epsilon is poorly matched to viral genomes, where mutation density can differ sharply between immunodominant surface regions and conserved domains; in SARS-CoV-2, default DBSCAN merged 99.76\% of points into one cluster, and tuned parameters ($\varepsilon = 0.15$, MinPts $= 5$) still yielded only 4 broad clusters~\cite{youn2025mutclust}.
HDBSCAN~\cite{campello2013hdbscan} and OPTICS~\cite{ankerst1999optics} address variable density through hierarchical stability or reachability plots, but remain generic spatial algorithms.
This limitation motivated MutClust~\cite{youn2025mutclust}, which replaces a global epsilon with an H-score-scaled, position-specific radius so biologically important positions receive wider neighborhoods.

\subsection{Deep Mutational Scanning as Ground Truth}
\label{subsec:bg_dms}

\textbf{Deep Mutational Scanning (DMS)} is an experimental technique that systematically measures the functional effects of all possible single amino acid substitutions within a protein~\cite{fowler2014dms}.
In a typical DMS experiment, a library of mutant variants---each carrying a single amino acid substitution at one position---is generated through saturation mutagenesis.
These variants are subjected to a selection assay (e.g., receptor binding, viral entry, or replication), and the relative fitness effect of each mutation is quantified by comparing variant frequencies before and after selection using deep sequencing.
The result is a comprehensive \textit{fitness landscape}: a matrix of fitness scores for every possible substitution at every position in the protein.

DMS is the gold standard for characterizing mutational effects because it is \textbf{unbiased} (every position and substitution is tested, not just those observed in nature), \textbf{comprehensive} (thousands of mutations are measured simultaneously in a single experiment), and \textbf{experimental} (fitness effects are directly measured rather than computationally predicted).
Unlike natural sequence variation, which reflects the combined effects of selection, drift, and sampling bias, DMS directly measures the intrinsic functional consequence of each mutation.

Several landmark DMS studies have shaped our understanding of viral protein fitness landscapes.
Starr et al.~\cite{starr2020dms} measured ACE2 binding affinity and protein expression for all single amino acid substitutions in the SARS-CoV-2 receptor-binding domain (RBD), revealing that most RBD positions are highly constrained for binding.
Dadonaite et al.~\cite{dadonaite2023dms} extended this to the full-length Spike protein by measuring pseudovirus cell entry fitness, identifying positions critical for viral entry beyond the RBD.
Lee et al.~\cite{lee2018h3n2dms} performed DMS of influenza H3N2 hemagglutinin (HA), demonstrating that DMS fitness measurements can predict the evolutionary fates of naturally circulating variants.
Bloom and Neher~\cite{bloom2023fitness} extended this further by inferring per-protein fitness effects across the entire SARS-CoV-2 proteome from public sequence phylogenies---rather than from direct DMS assays, since comprehensive DMS had then been performed on only two SARS-CoV-2 proteins---establishing a sequence-based reference complementary to direct DMS measurements.

A parallel lineage of \textbf{antibody-escape DMS} maps not viral fitness in isolation but the per-position escape from neutralization by specific antibodies or convalescent/vaccinee sera. Greaney et al.~\cite{greaney2021RBD} pioneered this design for the SARS-CoV-2 receptor-binding domain by measuring escape from a panel of monoclonal antibodies, and Greaney et al.~\cite{greaney2022spike} extended it to the full Spike with polyclonal sera (both SARS-CoV-2-specific). These SARS-CoV-2 antibody-escape DMS datasets, together with the Bloom-lab broader escape-DMS catalog, are the direct experimental antecedent of the \textbf{vaccine-escape cross-validation analysis} reported in Chapter~\ref{ch:discussion} (Section~\ref{sec:cross_validation_escape}; results tabulated in Chapter~\ref{ch:results} Table~\ref{tab:vaccine_escape_stage3}), which audits per-pathogen MutBench top-$k$ output against curated antibody-/vaccine-escape position lists for HIV-1 (Moore--Williamson 2015), H3N2 (Lee et al.~2018) and SARS-CoV-2. Layer~C ground truth in MutBench remains \emph{fitness}-DMS based and is sourced from Dadonaite et al.~2024 (SARS-CoV-2 cell entry), Lee et al.~2018 (H3N2 replicative fitness), Haddox et al.~2018 (HIV-1 viral growth in TZM-bl), Simonich et al.\ (RSV), Aditham et al.\ (Rabies) and Bakhache et al.\ (EV-A71); see Section~\ref{subsec:bg_dms} of this chapter and Table~\ref{tab:gt_positions} of Chapter~\ref{ch:mutbench}. The two ground-truth tracks are therefore orthogonal: antibody-escape DMS feeds the vaccine-escape audit, fitness-DMS feeds Layer~C, and each contributes a distinct per-position evaluation signal that is largely independent of literature-curated adaptive positives (Layer~A) and of purifying-selection negatives (Layer~B).

Despite its power, DMS has important limitations.
First, it is \textbf{expensive and labor-intensive}: each DMS experiment requires constructing a complete mutant library and performing selection assays, limiting its application to a small number of viral proteins.
Second, many surface-protein DMS datasets (especially for SARS-CoV-2 Spike and RSV F) use \textbf{pseudovirus systems} rather than authentic replication-competent virus, potentially missing fitness effects that depend on the complete viral replication cycle, although authentic replicative-fitness assays are also used in this dissertation's panel (e.g., MDCK culture for H3N2, RD culture for EV-A71).
Third, DMS measures fitness \textbf{in vitro}, which may not fully capture in vivo selective pressures such as immune evasion in the context of population-level immunity.
Finally, DMS coverage across the viral world remains sparse: the EVEREST v3 catalog~\cite{gurev2026everest_v3} covers only 45 viral DMS datasets even after explicitly targeting WHO priority viruses, and the catalog spans far fewer distinct viral proteins than the 40-virus WHO panel itself, leaving most pathogens (and most non-surface proteins within covered pathogens) without DMS measurements.
This availability gap means that DMS alone cannot serve as a universal ground truth for all pathogens---motivating MutBench's three-layer ground truth design (Chapter~\ref{ch:mutbench}), where DMS serves as Layer~C for the subset of pathogens with available data.

DMS data provides a natural ground truth for hotspot benchmarking: positions where mutations exhibit large fitness effects correspond to biologically functional sites that detection methods should identify~\cite{livesey2020using}.
Specifically, positions where the average DMS fitness effect of mutations is strongly negative (indicating functional constraint) or where specific mutations show large positive effects (indicating adaptive potential) are expected to overlap with biologically meaningful hotspots.
Livesey and Marsh~\cite{livesey2020using} demonstrated that DMS datasets can serve as unbiased benchmarks for computational variant effect predictors, establishing a precedent for the use of DMS as ground truth in method evaluation.

As a complementary approach, deep generative models have been shown to predict the effects of mutations from sequence variation alone~\cite{riesselman2018deep}, learning the statistical patterns of amino acid co-occurrence in evolutionarily related sequences to infer functional constraints without experimental data.
This approach was extended by EVE~\cite{frazer2021eve}, a Bayesian variational autoencoder (VAE) that learns evolutionary constraints from multiple sequence alignments (MSAs) without requiring experimental labels.
EVE achieves state-of-the-art performance on clinical variant classification for human proteins and provides uncertainty estimates through its Bayesian formulation.
The success of EVE and related generative models suggests that evolutionary information encoded in natural sequence variation captures much of the same functional signal as DMS experiments, albeit with lower resolution for individual mutations.


\subsection{Protein Language Models for Variant Effect Prediction}
\label{subsec:bg_plm}

Protein language models score mutations from evolutionary plausibility learned over large sequence corpora, offering structure- and assay-free features for variant-effect prediction.
ESM-2~\cite{lin2023esm2,rives2021biological} uses masked log-likelihood ratios, whereas Tranception~\cite{notin2022tranception} is autoregressive and can retrieve homologous context; both are relevant to MutBench because viral proteins may be underrepresented in cellular-protein-heavy training corpora.
Broader PLM and structure-aware systems include ProtTrans~\cite{elnaggar2022prottrans}, ESM-3~\cite{hayes2024esm3}, ESM-IF~\cite{hsu2022esmif}, SaProt~\cite{su2024saprot}, and ProteinMPNN~\cite{dauparas2022proteinmpnn}; ESM-3 is excluded from the headline 8{,}580-evaluation grid because its Cambrian non-commercial license conflicts with the MIT-redistribution requirement (Chapter~\ref{ch:mutbench}, Section~\ref{para:external_licenses}), and structure-aware channels are deferred as future work.
EVEscape~\cite{thadani2023evescape} shows that generative fitness, accessibility, and residue dissimilarity can be productively combined; MutBench adopts only an EVEscape-style composite that substitutes ESM-2 mutation surprise for the EVE term (Section~\ref{sec:scoring_detection}).
AlphaMissense~\cite{cheng2023alphamissense} is important but human-variant-centered, while EVEREST~\cite{gurev2025everest,gurev2026everest_v3} reports PLM underperformance on viral proteins as a cross-study parallel, not independent confirmation of MutBench's pathogen-dependence result (Section~\ref{subsec:bg_benchmarking_gap}).
ProteinGym~\cite{notin2023proteingym} ranks predictors for per-mutation fitness regression over DMS assays; MutBench instead asks which information types, including ESM-2 masked-marginal and semantic-change scores, identify mutation-enriched regions across pathogens (Section~\ref{sec:information_analysis}).

\subsection{Viral Genomic Surveillance Platforms}
\label{subsec:bg_surveillance}

Real-time viral genomic surveillance platforms provide the operational context in which hotspot detection methods are deployed, but serve a fundamentally different purpose.

\textbf{Nextstrain}~\cite{nextstrain2018hadfield} integrates phylogenetic analysis with geographic and temporal metadata to track pathogen evolution in real time. Its primary function is to answer ``which lineages are currently spreading and where?''---for example, tracking the geographic expansion of SARS-CoV-2 Omicron sublineages or seasonal influenza clades. Nextstrain visualizes the evolutionary relationships among circulating variants but does not evaluate or compare methods for detecting mutation-prone positions.

\textbf{PANGO}~\cite{rambaut2020pangolin} provides a standardized lineage classification system for SARS-CoV-2, assigning labels (e.g., BA.2.86, JN.1) to sequences based on their phylogenetic placement. Its purpose is taxonomic organization---labeling ``what is this variant?''---rather than identifying which genomic positions are likely to accumulate future mutations of functional significance.

These platforms cannot substitute for hotspot detection benchmarking for three reasons:
(1)~they operate primarily at the \textit{lineage/variant level} (which variant is spreading), not at the \textit{position level} (which genomic position is mutation-prone);
(2)~they are predominantly \textit{retrospective} (tracking variants that have already emerged), although the Bedford-lab forecasting line (Łuksza \& Lässig~\cite{luksza2014predictive}; Huddleston et al.~\cite{huddleston2020integrating}) extends Nextstrain toward short-horizon clade-level prediction at the \emph{lineage} rather than per-position level;
and (3)~they do not compare or evaluate detection algorithms---they are end-user tools, not method evaluation frameworks.
The distinction is therefore one of \emph{primary function} rather than absolute exclusion: surveillance answers ``what has happened, and which clades are growing?'' while hotspot detection aims to answer ``where, at the residue level, might important mutations occur next?''
MutBench addresses the latter question by providing a systematic framework for evaluating which detection methods most effectively identify biologically meaningful mutation-prone positions.

Other surveillance-related resources include \textbf{GISAID}~\cite{shu2017gisaid}, the primary repository for pathogen genome sequences providing the raw data from which mutation frequencies are computed, and \textbf{outbreak.info}, which aggregates variant prevalence data from GISAID.
These databases are essential data \textit{sources} for hotspot detection studies (including the present dissertation), but they do not themselves provide detection or benchmarking capabilities.
The relationship between surveillance platforms and hotspot detection is thus complementary: surveillance platforms provide the data and operational context, while benchmarking frameworks like MutBench provide the methodological infrastructure for evaluating detection methods applied to that data.

\subsection{Feature Importance Analysis in Viral Mutation Prediction}
\label{subsec:bg_feature_importance}

Prior feature-importance studies motivate MutBench's multi-source analysis but remain mostly single-pathogen: Maher et al.~\cite{maher2022predicting} found epidemiological features dominated sequence features for SARS-CoV-2 variant emergence, Rodriguez-Rivas et al.~\cite{rodriguez2022epistatic} showed epistatic DCA models improved escape prediction (AUC 0.76 vs.\ 0.63), Hie et al.~\cite{hie2021learning} linked PLM semantic change to antigenic novelty, and EVEscape~\cite{thadani2023evescape} combined fitness, accessibility, and residue dissimilarity.
Coupling-aware methods such as DCA, EVcouplings, and GREMLIN remain outside MutBench's 20-scoring panel (Section~\ref{sec:scoring_detection}), so the scoring~$\times$~pathogen interaction in Chapter~\ref{ch:results} is an information-type contrast within a single-position regime, not an upper bound on epistasis-aware methods.
MutBench addresses the remaining gap by evaluating 20 scoring channels from 10 biological information families across 11 pathogens and reporting pathogen-specific feature recommendations (Section~\ref{sec:information_analysis}).

\subsection{Algorithm Selection and Benchmarking Methodology}
\label{subsec:bg_algorithm_selection}

The finding that the best-performing detection method depends on the pathogen (Chapter~\ref{ch:results}) is an instance of Rice's algorithm selection problem~\cite{rice1976algorithm}.
Rice formalized this as a mapping from a \textit{problem space} (characterized by measurable features of each problem instance) through a \textit{feature space} to an \textit{algorithm space}, with the goal of selecting the algorithm that maximizes a given \textit{performance measure} for each instance.
In the context of hotspot detection, each pathogen constitutes a problem instance characterized by features such as genome size, evolutionary rate, and mutation density distribution; the algorithm space comprises the set of scoring--detection combinations; and the performance measure is benchmark-stage-specific---hotspot-score (recall+precision+stability) for Stage~1 SARS-CoV-2 validation and MCC for the Stage~2 11-pathogen comparison (see Section~\ref{subsec:9pathogen_metrics}).

This perspective is closely related to the \textbf{No Free Lunch (NFL) theorem}~\cite{wolpert1997nfl}, which states that no single optimization algorithm outperforms all others across all possible problem instances.
Importantly, NFL does \textit{not} imply that all algorithms perform equally on any given problem---rather, it implies that superior performance on one class of problems necessarily comes at the cost of inferior performance on others.
For hotspot detection, this means that a scoring type optimized for one pathogen's evolutionary characteristics may be suboptimal for another, motivating the systematic cross-pathogen comparison that MutBench provides.

This framework has been formalized in the machine learning community through meta-learning~\cite{kerschke2019algorithm,bischl2016aslib} and automated algorithm selection (AutoML)~\cite{feurer2015automl,kotthoff2016algorithm}, where the goal is to automatically select the best-performing model or pipeline configuration for a given dataset.
The parallel to viral hotspot detection is direct in principle: pathogen characteristics (genome size, mutation rate, population structure, selective pressure regime) constitute the problem-instance features, and scoring--detection pipelines constitute the algorithm space. The concrete pathogen-to-scoring recommendation produced by MutBench under this framing is presented in Chapter~\ref{ch:discussion} as an application of the prior-work framework, not as part of the background itself.
However, the algorithm selection framework has not previously been applied to viral genomics in this form, where the ``problem instances'' are pathogens with diverse evolutionary characteristics and the ``algorithms'' are scoring--detection pipelines.

The benchmarking methodology adopted in this dissertation follows established principles from computational biology.
Weber et al.~\cite{weber2019benchmarking} proposed guidelines for neutral benchmarking, articulating three core requirements: (1)~\textit{independent evaluation}, where the benchmark designer is distinct from the method developers; (2)~\textit{multiple metrics}, to capture different aspects of performance and avoid optimization for a single criterion; and (3)~\textit{controlled experimental design}, ensuring that all methods are evaluated on identical data under identical conditions. Mangul et al.~\cite{mangul2019systematic} extended these principles to a community-wide audit of omics computational tools, documenting systematic disparities in tool reliability that motivated the present study's emphasis on multi-pathogen rather than single-pathogen evaluation.
Boulesteix et al.~\cite{boulesteix2017towards} further emphasized the importance of avoiding ``self-assessment bias'' where method developers evaluate their own methods on self-selected datasets, which systematically inflates reported performance.
MutBench adheres to these principles---with the partial exception that MutClust~\cite{youn2025mutclust}, developed in the same lab as the present dissertation, is one of the 14 evaluated detection families and therefore violates Weber's strict ``benchmark designer $\neq$ method developer'' criterion. This self-assessment-bias risk (in the sense of Boulesteix et al.~\cite{boulesteix2017towards}) is mitigated, but not eliminated, by three design choices: (i)~all 14 detection families are evaluated under identical preprocessing, scoring, and metric pipelines; (ii)~the three-layer ground truth (literature-curated adaptive positives, purifying-selection negatives, DMS-anchored Layer~C) is constructed independently of any method's predictions; and (iii)~MutClust's hyperparameters were frozen prior to Stage~2 sweep execution. The remaining design influence is acknowledged as a limitation in Chapter~\ref{ch:discussion}. With this caveat, MutBench evaluates all methods using a composite metric (Stage~1 hotspot-score: recall+precision+stability) for in-depth SARS-CoV-2 validation and MCC for the cross-pathogen Stage~2 comparison (see Section~\ref{subsec:9pathogen_metrics} for the rationale).

\subsection{Benchmarking Gap in Viral Hotspot Detection}
\label{subsec:bg_benchmarking_gap}

The absence of benchmark studies for hotspot detection methods in viral genomics stands in contrast to cancer genomics, where Bailey et al.~\cite{bailey2018comprehensive} conducted a comprehensive PanCancer analysis applying 26 computational tools to 9,423 tumor exomes under uniform conditions, producing a widely used catalog of cancer driver genes.
No analogous \emph{detection-task} benchmark exists for viral mutation hotspots: while several recent fitness-regression benchmarks now span comparable virus panels, the region-level hotspot-detection task itself remains unstandardized---detection methods are evaluated using ad hoc validation against small sets of known functional mutations, without standardized datasets, metrics, or ground truth definitions.
Related efforts in adjacent domains include:
(1)~ProteinGym~\cite{notin2023proteingym,notin2025proteingym_v13}, which benchmarks variant effect predictors using 217 DMS assays and 90+ baseline models (as of v1.3) but addresses a per-substitution fitness-regression task rather than hotspot region detection;
(1b)~\textbf{ViroGym}~\cite{virogym2026} (Mar 2026), which benchmarks PLM and alignment-based scorers across 13 viruses, 79 DMS assays, 552{,}937 sequences, and 7 phenotypic categories including immune escape---larger virus count than MutBench, but focused on per-variant fitness rather than per-region detection;
(2)~ConDor~\cite{morel2024condor}, which detects convergent mutations across phylogenies but is a single tool rather than a benchmark framework;
(3)~Pons et al.~\cite{pons2020computational}, who reviewed 40+ cancer hotspot detection methods but produced a descriptive survey without a unified benchmark, and cancer-specific models do not transfer to viral evolutionary dynamics;
and (4)~individual SARS-CoV-2 studies using entropy~\cite{harvey2021tracking} or clustering that validate against small variant sets without standardized ground truth.
Among the 35+ mutation clustering tools catalogued in the computational oncology literature~\cite{pons2020computational}, none was designed for or validated on viral population-level mutation data. Cancer hotspot tools (OncodriveCLUST~\cite{tamborero2013oncodriveclust}, HotMAPS, MutSpot) require tumor cohort somatic mutation calls and human-genome-specific background models (trinucleotide signatures, replication timing). Phylogenetic selection tools (MEME, FUBAR, FEL, SLAC, BUSTED) detect per-codon dN/dS signals but do not perform spatial clustering of hotspot regions. The only virus-specific hotspot clustering tool is MutClust~\cite{youn2025mutclust}, validated on a single pathogen (SARS-CoV-2). This gap---the absence of systematic multi-pathogen hotspot detection benchmarking---is what MutBench addresses.
Thus, a systematic benchmark with multi-source ground truth, composite evaluation metrics, and cross-pathogen validation remains absent for viral hotspot detection.

\subsubsection*{Variant Effect Prediction Benchmarks}

Recent years have seen significant progress in benchmarking variant effect prediction (VEP) methods on proteins. ProteinGym~\cite{notin2023proteingym} established a large-scale benchmark comprising 217 DMS assays and 2.7 million substitution variants, primarily for human and model organism proteins. EVE~\cite{frazer2021eve}, a Bayesian VAE trained on evolutionary sequences, demonstrated that unsupervised generative models can predict disease-associated variants without clinical labels. EVEscape~\cite{thadani2023evescape} extended this framework to viral immune escape prediction by combining EVE fitness scores with structural accessibility and physicochemical dissimilarity, successfully forecasting SARS-CoV-2 Omicron escape mutations from pre-pandemic data.

Most recently, EVEREST~\cite{gurev2025everest,gurev2026everest_v3} specifically benchmarked VEP methods across WHO priority viruses; the v3 expansion (Jan 2026) covers 45 viral DMS datasets, 31 forecasting clades across 4 viruses, and reliability assessment across 40 WHO priority RNA viruses spanning 16 viral families. EVEREST finds that protein language models substantially underperform alignment-based approaches on viral proteins, that reliability fails for over half of the WHO-40 panel (especially surface glycoproteins), and that prediction reliability varies dramatically across pathogens. Because MutBench's LOPO 0/11 match rate is itself null-consistent under permutation (Chapter~\ref{ch:results}), we describe this as a cross-study \emph{parallel} on pathogen-dependence rather than as independent corroboration of non-transferability per se.
The PLANT framework~\cite{plant2025antigenic} and Shah et al.~\cite{shah2024h3n2ml} provide PLM-based and ML-based H3N2 antigenic cartography that explicitly competes with WHO vaccine-strain selection, building on the foundational antigenic-cartography paradigm of Smith et al.~\cite{smith2004antigenic}; MutBench's H3N2 self-consistency check (Section~\ref{sec:cross_validation_escape}) is positioned as a detection-side complement to these prediction-side cartographies, not as a competitor.
A particularly relevant finding from EVEREST is that the performance gap between methods varies by viral family: alignment-based methods such as EVE consistently outperform PLMs across most viral proteins, but the magnitude of this advantage differs substantially between, for example, influenza HA and SARS-CoV-2 Spike.
This pathogen-dependent performance variation parallels the central finding of MutBench (Chapter~\ref{ch:results}).
AlphaMissense~\cite{cheng2023alphamissense} demonstrated proteome-wide variant classification using an AlphaFold-derived architecture, though its focus on human proteins limits direct viral applicability.

MutBench's unique contribution is the \textit{scoring $\times$ pathogen interaction analysis}: while EVEREST establishes that method performance varies across viruses, MutBench quantifies this interaction through three-way ANOVA ($\omega^2 = 0.296$) and demonstrates that the choice of \textit{biological information source} (not just the computational method) is the dominant performance factor. This finding---that the optimal information source is pathogen-dependent within the 20-type/39-variant design, with the interaction as the largest modeled variance component---has no counterpart in EVEREST or ProteinGym, which treat the scoring/feature set as fixed.
Specifically, MutBench provides three capabilities absent in EVEREST and other VEP benchmarks:
(1)~quantification of the scoring--pathogen interaction through three-way ANOVA, revealing that the scoring $\times$ pathogen interaction ($\omega^2 = 0.296$; 11-pathogen cluster-bootstrap 95\% interval [0.201, 0.346]) is the largest modeled variance component (within-cell residual $\omega^2 \approx 0.39$);
(2)~cross-validation via vaccine escape enrichment for the 3 of 11 pathogens with curated escape annotations, with HIV-1 (7.19$\times$, Bonferroni-corrected $p = 2.5 \times 10^{-16}$; escape set 82\% disjoint from Layer~A) as the primary external anchor and H3N2 9.36$\times$ reported as a self-consistency check given 69\% Layer~A overlap (circularity audit in Table~\ref{tab:vaccine_escape_stage3});
and (3)~a concrete pathogen-to-scoring recommendation table enabling application to new pathogens without requiring DMS experimental data.
Furthermore, MutBench introduces the vaccine escape cross-validation paradigm with a per-pathogen circularity audit: HIV-1 provides the cleanest external validation, while H3N2 and SARS-CoV-2 contribute self-consistency and exploratory evidence respectively---a distinction absent from EVEREST or ProteinGym.
Table~\ref{tab:benchmark_comparison} summarizes the key differences between MutBench and these related benchmarks.

\begin{table}[htbp]
\centering
\caption[Comparison of MutBench with related benchmarks]{Cross-benchmark comparison: task, scope, ground-truth, methods compared. The ViroGym row reflects the March 2026 arXiv preprint (pre-peer-review) and is included for completeness; EVEREST v3 (Jan 2026) is a bioRxiv preprint and ProteinGym v1.3 is a versioned release/update---both are well-established community benchmarks but the v3 / v1.3 versions specifically have not been confirmed peer-reviewed in this dissertation's bibliography.}
\label{tab:benchmark_comparison}
\footnotesize
\begin{adjustbox}{max width=\textwidth}
\begin{tabular}{lccccc}
\toprule
& \textbf{MutBench} & \textbf{ViroGym} & \textbf{EVEREST v3} & \textbf{ProteinGym v1.3} & \textbf{Bailey 2018} \\
\midrule
Domain & Viral hotspot & Viral VEP & Viral VEP & General VEP & Cancer driver \\
Task & Detection & Prediction & Prediction & Prediction & Detection \\
Pathogens/proteins & 11 viruses & 13 viruses (79 DMS) & 4--40 viruses (45 DMS) & 217 proteins & Cancer genes \\
Ground truth & 3-layer & DMS+escape & DMS & DMS & Patient cohorts \\
Methods compared & 780 combos & $\sim$10 PLM/MSA & $\sim$20 VEP & 90+ VEP & 26 tools \\
Cross-pathogen ANOVA & Yes ($\omega^2$) & No & No & No & No \\
\bottomrule
\end{tabular}
\end{adjustbox}
\par\vspace{0.5em}
\footnotesize Direct performance comparison between MutBench and ProteinGym, EVEREST, or ViroGym is not possible because the tasks differ fundamentally: VEP benchmarks (including ViroGym's per-variant fitness/escape regression and forecasting tasks) evaluate Spearman correlation or AUROC between predicted and measured per-variant outcomes, while MutBench evaluates MCC for binary per-region hotspot classification. The benchmarks are complementary rather than competing.
\end{table}

This gap---systematic multi-pathogen benchmarking for hotspot detection---is the fundamental motivation for MutBench (Chapter~\ref{ch:mutbench}), the core study of this dissertation.



%% ===================================================================
%%  Related Research Summary
%% ===================================================================
\vspace{1em}
\noindent\textbf{Related research summary.}
This section has reviewed the two bodies of prior work that directly inform MutBench.
Section~\ref{sec:bg_hotspot} surveyed viral mutation hotspot detection methods---from frequency-based and entropy-based scoring through density-based clustering to protein language models and selection-based approaches---and identified a critical benchmarking gap: no standardized framework exists for evaluating and comparing these methods across multiple pathogens with independent ground truth.
The only virus-specific hotspot detection tool, MutClust~\cite{youn2025mutclust}, relies on a single pathogen, single information type (H-score), and self-referential validation, leaving open the question of how different scoring approaches and detection algorithms compare across diverse viral pathogens.
Together, these gaps motivate the design of MutBench as a systematic, multi-pathogen benchmark with three-layer ground truth, composite evaluation metrics, and statistical analysis of scoring--pathogen interactions.
Coupling-aware methods (DCA, EVcouplings, GREMLIN, and similar pairwise/epistatic models) remain explicitly out of scope of the 20-scoring panel; the scoring~$\times$~pathogen interaction $\omega^2 = 0.296$ and all downstream Stage~2/3 conclusions should therefore be read within a single-position scoring regime, with epistasis-driven hotspot detection deferred as future work (Chapter~\ref{ch:discussion}, Section~\ref{sec:future_work}).
Chapter~\ref{ch:mutbench} presents the MutBench framework in detail, including the construction of multi-layer ground truth datasets for 11 pathogens, the benchmark design comprising 20 scoring types across 6 categories and 14 detection method families, and the factorial ANOVA methodology for quantifying scoring--pathogen interactions.
